% ==============================================================================
% РАЗДЕЛ 2.4 (DISS-005-B1445 .. DISS-005-B1469)
% Выводы по главе 2
% Замечания: DISS-005-C0100 .. DISS-005-C0102 (3 шт.)
% ==============================================================================

% DISS-005-B1445
\section{Выводы}
\label{sec:ch02_conclusions}

% DISS-005-B1446
Проведенные\commentnote{DISS-005-C0100}{Алексей Евдаков [2]}{Тоже накидано нейронкой. Пока что сойдёт} в настоящей главе исследования и анализ современных подходов к модернизации устройств быстродействующего автоматического ввода резерва (БАВР) позволяют сделать следующие основные выводы:\par\vspace{0.2cm}

\begin{enumerate}[label=\arabic*)]
    % DISS-005-B1447
    \item Применение катушек Роговского (КР) в качестве первичных преобразователей тока для реле направления мощности (РНМ) в устройствах БАВР является перспективным решением, позволяющим существенно улучшить динамические характеристики и надежность измерительных органов.
    \begin{itemize}
        % DISS-005-B1448
        \item Катушки Роговского обеспечивают более быстрый переход в установившийся режим измерения по сравнению с традиционными ТТ, особенно в условиях наличия значительной апериодической составляющей в токе короткого замыкания. Время установления первой гармоники для КР обычно не превышает периода промышленной частоты (20~мс для отечественных сетей).
        % DISS-005-B1449
        \item Отсутствие ферромагнитного сердечника исключает проблему насыщения КР при больших токах, что критически важно для корректной работы защит в тяжелых аварийных режимах.
        % DISS-005-B1450
        \item Катушки Роговского менее подвержены влиянию апериодической составляющей, что повышает точность расчета первой гармоники и стабильность работы РНМ.
        % DISS-005-B1451
        \item Исследования влияния различных факторов, таких как параметры дискретизации, оконные функции БПФ и наличие высших гармоник, показали возможность достижения высокой точности измерений с КР при правильном выборе алгоритмов обработки сигнала.
        % DISS-005-B1452
        \item Практическое внедрение терминалов БАВР с КР подтвердило их эффективность и надежность в реальных условиях эксплуатации. Перспективным является также гибридный подход, сочетающий сигналы от ТТ и КР для оптимизации работы защит в различных режимах.
    \end{itemize}
    % DISS-005-B1453
    \item Совершенствование алгоритмов функционирования восстановления нормального режима (ВНР) является необходимым условием для обеспечения надежной работы БАВР в сложных и нештатных ситуациях.
    \begin{itemize}
        % DISS-005-B1454
        \item Внедрение блокировок ВНР при наличии однофазных замыканий на землю, при превышении допустимого угла между напряжениями источников, а также по разности частот позволяет предотвратить некорректные переключения, способные привести к развитию аварии (например, междуфазному КЗ при включении на ОЗЗ или к несинхронному включению).
        % DISS-005-B1455
        \item Использование полноценного контроля синхронизма (ANSI 25) перед параллельным включением источников значительно повышает надежность ВНР, особенно в сетях с собственной генерацией или при работе с удаленными источниками.
        % DISS-005-B1456
        \item Введение временной селективности при срабатывании нескольких устройств ВНР на одном энергообъекте позволяет снизить суммарные броски тока и перетоки мощности, предотвращая перегрузку оборудования и возможные каскадные отключения.
        % DISS-005-B1457
        \item Предложенные модернизации алгоритмов ВНР способствуют минимизации рисков ложных или излишних срабатываний автоматики, повышая общую устойчивость электроснабжения.
    \end{itemize}
\end{enumerate}

% DISS-005-B1458
Внедрение рассмотренных решений в практику проектирования и эксплуатации систем РЗА способствует повышению стабильности электроснабжения промышленных предприятий и других ответственных потребителей, минимизируя экономический ущерб от перерывов в питании.\par\vspace{0.2cm}

% DISS-005-B1459
Комплексное рассмотрение вопросов применения современных измерительных технологий на базе катушек Роговского и усовершенствованных алгоритмов восстановления нормального режима показывает, что повышение общей надежности и эффективности функционирования БАВР достигается как за счет улучшения качества получаемой измерительной информации, так и за счет более продуманной и современной логики функционирования данной автоматики.\par\vspace{0.2cm}

% DISS-005-B1460
Данным модернизациям\commentnote{DISS-005-C0101}{Алексей Евдаков [2]}{По мере развития / написания глав нужно будет поправить.}, основанным на современных подходах обработки и анализа больших данных будут посвящены следующие главы. В них будет описан весь процесс сбора и подготовки данных (глава ХХ). Предложены варианты статистического анализа реальных осциллограмм для выявления и проверки различных теорий (раздел глав ХХХ). И наконец, рассмотрены варианты применения нейронных сетей для модернизации алгоритмов БАВР (раздел глав ХХХ).\par\vspace{0.2cm}

% DISS-005-B1461

% DISS-005-B1462
\paragraph{На подумать}\commentnote{DISS-005-C0102}{Алексей Евдаков [2]}{Прост пометки, чтобы не забыть.}

% DISS-005-B1463
Статистический анализ, что стоит и можно сделать. Пишу чтобы не забыть\par\vspace{0.2cm}

\begin{itemize}[leftmargin=*,itemsep=0pt]
    % DISS-005-B1464
    \item анализ времён используемых выключателей (данное уже подготовил, так почему бы не прописать)
    % DISS-005-B1465
    \item анализ осциллограмм с ОЗЗ и оценка реальных перенапряжений. (хочу сделать статью совместно с Шуиным, уже в процессе обработки)
\end{itemize}

% DISS-005-B1466

% DISS-005-B1467
\paragraph{Данные.}

% DISS-005-B1468
Сейчас у меня будет большой разбор анализа реальных осциллограмм. Но надо бы собрать и симулированные данные, их очень не хватает, и они нужны будут.\par\vspace{0.2cm}

% DISS-005-B1469

